\section{Capturing a cluster and paying for a cutoff match}
\label{sec:ope-cutoff}
The local match used in Section~\ref{sec:ope-grid} is constructed before
any continuum limit. One follows a product of local observables down
through the scales, merging groups whenever their parent sites cease
to be separated. The bookkeeping must be independent of the number of
replicated cells. Otherwise it would be unusable in the centered root
estimate. We give the geometric selection and the scale accounting
explicitly, and state where the remaining local RG estimates enter.

\subsection{Regrouping without a growing arity cost}
Let a cluster have $h\geq1$ marked sites and distinguish one of its
marks. At a given scale, call sites separated if their periodic
maximum-coordinate distance is at least $B^2$, where $B\geq1$ is the
purchased support radius. Among the subsets containing the distinguished
mark choose one of maximal cardinality with separated sites. Every
omitted mark lies at distance less than $B^2$ from a chosen mark:
otherwise it could be added. Assign each omitted mark to one such
representative, and assign each chosen mark to itself. This defines a
partition with a chosen representative in every part. If the original
sites were not separated, the number of parts strictly decreases.

Copies in different centered cells have the same relative geometry,
up to translations and signed coordinate transformations, which preserve
distance. Use the same partition in each copy. Cross-cell separation
and separation of the chosen representatives in the home cell then
give separation of all representatives. In particular the number of
regroupings is at most $h-1$, regardless of how many cells are copied.

The cost of connecting a mark to its representative can be charged
one scale before their parents meet. To check the geometry, let the
block ratio be $L$, let $K$ be the parent's support radius and $B$ the
child's, and let $x,y$ be two child sites whose parents have distance
less than $K^2$. Each child is at distance less than $L/2$ from the
embedded parent. The triangle inequality gives
\[
 \operatorname{dist}(x,y)<L(K^2+1)\leq L(K+B+1)^2=:c.
\]
A pair at distance at most $cB$ can be covered by at most $(2c+1)^4$
$B$-blocks: both points lie in the radius-$cB$ halo of one point,
which is tiled by at most $2c+1$ blocks per coordinate. Since $B\geq1$,
the preceding pair therefore has covering number at most
\begin{equation}\label{eq:ope-pair-halo}
 (2c+1)^4\leq16(1+c)^4.
\end{equation}
This is the halo charged to the individual grouping operation. It
does not assert that the entire cluster lies in a single small halo
at every intervening scale.

Here are the two local RG estimates used in the induction. A group
has a positive weight $W$ measuring the allowed size of its marked
representation. If its parents stay separated, descent through a block
of ratio $L$ changes its weight to $W/L^4$. If $p$ groups merge into
$q<p$ groups at a scale with connecting-cover bound $\ell$, their
new weights can be taken to be
\begin{equation}\label{eq:ope-group-weight}
 \widetilde W_A=\frac18\prod_{r\in A}(\lambda W_r),\qquad
 \lambda=\frac{\delta^2}{h}
       e^{-(\nu+1/2)\ell}.
\end{equation}
Here $\nu$ is the purchased support exponent and $\delta>0$ the local
smallness parameter. These statements also retain the exact marked
polynomial after regrouping, its covariance between copies, its
integrability, and its permitted local representation. The analytic
proofs of these two local representation estimates remain construction
inputs. In particular their weight rules alone are not a proof that
a marked representation exists.

Choose one $0<\alpha\leq1$ no larger than every grouping factor
$\lambda$ that will be needed. If $j$ is the initial scale, $i$ the
current scale and $s\leq h$ the current number of groups, put
\[
 C(i,s)=\left(\frac{a_j}{a_i}\right)^{4h}
         (\alpha/8)^{h(h-s)}\prod_{r=1}^h W_r^{\mathrm{initial}}.
\]
We claim that the product of the current group weights is at least
$C(i,s)$. This is equality initially. For an ordinary descent,
\[
 C(i,s)=L_i^{-4h}C(i+1,s)
 \leq L_i^{-4s}\prod_r W_r,
\]
since $s\leq h$ and $L_i\geq1$.
For a genuine grouping, $q<p\leq h$ and $\alpha/8\leq1$ imply
\[
 C(i,q)\leq C(i,p)(\alpha/8)^h.
\]
Thus
\[
 8^q C(i,q)\leq\alpha^h C(i,p)
       \leq\lambda^p\prod_r W_r.
\]
We used $8^q\leq8^h$ and $\alpha^h\leq\alpha^p\leq\lambda^p$.
The last expression is exactly the product of the new numerators
in \eqref{eq:ope-group-weight}, because the groups partition the
old marks. This proves the induction after the subsequent descent
as well. The regrouping of marked polynomials commutes with descent:
it is a fixed linear map on a finite coefficient space, so it commutes
with integration against the transition kernel. Its composition law
is the composition law of the partitions. These observations keep the
induction on the original exact product.

The initial single-field estimates give weights
$W_r^{\mathrm{initial}}=\delta_j\eta_j p_r(j)/8$, where $p_r(j)>0$
is the paid weight of the $r$th field. At the final parent scale $m+1$
merge all remaining groups. Their number $s$ is at most $h$, and the
last grouping factor is also at least $\alpha$. The terminal merge
uses the weight $\lambda^s\prod_r W_r$ without the factor $1/8$
reserved at intermediate regroupings; it supplies the final local
projection directly. Since $\alpha^h\leq\lambda^s$ and
$C(m+1,0)\leq C(m+1,s)$, the resulting weight is at least
\begin{equation}\label{eq:ope-initial-weight}
 X=\alpha^h
 \left(\frac{a_j}{a_{m+1}}\right)^{4h}
 (\alpha/8)^{h^2}
 (\delta_j\eta_j/8)^h\prod_{r=1}^h p_r(j).
\end{equation}
The final local projection estimate produces a degree-$q$ coefficient
whose coefficient-size and remainder norm are each bounded by
\begin{equation}\label{eq:ope-terminal-coefficient}
 P=a_{m+1}^q\,\mathcal B_m\frac4X.
\end{equation}
Here $\mathcal B_m$ is the purchased local projection bound. This
projection estimate, together with the marked-representation estimates
above, is still an analytic input. The preceding argument proves the
precise scale cost it must be given; no factor depends on the number
of copies.

\subsection{A single subpower allowance for all intermediate levels}
The halo estimate ensures that each grouping factor has zero
logarithmic rate. Explicitly the purchased bound is
$(\nu_{i+1}+1)\ell_i\leq16\sqrt{\widehat\beta_i}$,
while $\sqrt{\widehat\beta_i}/s_i\to0$ and
$\log\delta_{i+1}=o(s_i)$. Taking logarithms of
\eqref{eq:ope-group-weight} therefore gives
$\log\lambda_i=o(s_i)$ for fixed $h$. These are schedule and
running-coupling estimates; their production is a separate input.

This pointwise-in-scale fact also gives a uniform allowance for all
previous scales. If $u_i>0$ and $\log u_i=o(s_i)$ for an increasing
$s_i\to\infty$, set $\alpha_j=\min(1,u_0,\ldots,u_j)$. Then
$\alpha_j\leq1$, $\alpha_j\leq u_i$ for $i\leq j$, and
$\log\alpha_j=o(s_j)$. Indeed, for any $\epsilon>0$, eventually
$u_i\geq e^{-\epsilon s_i}\geq e^{-\epsilon s_j}$ for $i\leq j$
in the tail. The positive minimum of the finite initial segment is
also eventually at least $e^{-\epsilon s_j}$. Hence
$e^{-\epsilon s_j}\leq\alpha_j\leq1$. Apply this to the grouping
factors. Arbitrarily many intermediate scales thus require only one
subpower allowance for the fixed cluster arity.

\subsection{The two output prices}
Take sequences of levels $k,m,j$ of logarithmic rates $b,a,c$ relative
to $t_n\to\infty$, with $0<b<a<c$. Here $k$ sets the centered test
scale, $m$ the matching scale, and $j$ the insertion grid. All level
sequences tend to infinity. Let the field degrees be $d_r$, put
$w=\sum_r d_r$, and define
\[
 \tau=cw+4h(c-a).
\]
The calibrated single-field weight has the form
$p_r(j)=a_j^{d_r}\mathcal W_r(j)$, with
$\log\mathcal W_r(j)=o(s_j)$. All factors in
\eqref{eq:ope-initial-weight} other than the displayed mesh powers
have zero logarithmic rate. The gentle schedule gives rate $a$ also
at $m+1$. Consequently $X$ has rate $-\tau$.

Set
\begin{equation}\label{eq:ope-terminal-price}
 E=4(1+\zeta_m)(2P),\qquad V=2E,
 \qquad z=\exp\bigl(((q+\tfrac1{16})(a-b)-\tau)t_n\bigr).
\end{equation}
Equations \eqref{eq:ope-terminal-coefficient} and
\eqref{eq:ope-terminal-price} give rate $\tau-qa$ for both $E,V$.
The final comparison of the local coefficient with the normalized
reference fields has the following quantitative form. Provided its
positive matching allowance $z$ satisfies
\begin{align}
 \mathcal M(z;k,m)&\leq\delta_m X/(16\mathcal B_m),
       \label{eq:ope-match-pay-one}\\
 32(1+\zeta_m)z\,a_k^{-q}
       E/G(m,k)&\leq\vartheta_q(k)/2,
       \label{eq:ope-match-pay-two}
\end{align}
it produces a reference label $Q_{N,p}$ with quota at most $V$ and
centered discrepancy price at most
\begin{equation}\label{eq:ope-cutoff-discrepancy}
 A=2\mathcal S(k)/z.
\end{equation}
Here $G(m,k)=\exp((s_m-s_k)/8)$, $\vartheta_q$ is positive with
zero logarithmic rate, and $\mathcal M$ is the allowance transported
from the centered scale to $m$. Its bound is
\[
 \mathcal M(z;k,m)\leq
 z(a_m/a_k)^{q+1}\frac{\vartheta_q(m)}{\vartheta_q(k)}.
\]
If the centered scale has first been opened to a finer admissible
level, use that opened level in these formulas. The opening retains
rate $b$, as checked below. The actual comparison producing
\eqref{eq:ope-cutoff-discrepancy}, including its representation and
reference-field identification, remains to be proved in the
construction part; these displayed inequalities specify exactly the
budget it needs.

Both payment conditions hold eventually by a strict comparison of
rates. The right side of \eqref{eq:ope-match-pay-one} has rate
$-\tau$, whereas the upper bound for its left side has rate
\[
 (q+\tfrac1{16})(a-b)-\tau-(q+1)(a-b)
       =-\tau-\tfrac{15}{16}(a-b)<-\tau.
\]
The left side of \eqref{eq:ope-match-pay-two} has rate
\[
 (q+\tfrac1{16})(a-b)-\tau+qb+(\tau-qa)
       -\tfrac18(a-b)=-\tfrac1{16}(a-b)<0,
\]
while its positive right side has rate zero. If positive sequences
have unequal logarithmic rates, division and taking logarithms shows
that the one with the smaller rate is eventually smaller. This
justifies both comparisons, including all their fixed constants.
Since $\mathcal S$ has upper rate zero, \eqref{eq:ope-cutoff-discrepancy}
has upper rate $\tau-(q+1/16)(a-b)$, as required in
Input~\ref{input:ope-match}.

Here is the opening check. Write $K(k)\geq k$ for the opened level,
$u=a_k/a_{K(k)}\geq1$, and $B_{K(k)}\geq1$ for its positive vacuum
budget. The schedule ensures, for every $\epsilon>0$,
\[
 u^4 B_{K(k)}\leq\epsilon s_{K(k)}\quad\hbox{eventually in }k.
\]
Since $s_{K(k)}-s_k=\log u\leq u\leq u^4B_{K(k)}$,
we have $s_k/s_{K(k)}\to1$. Opening a sequence of rate $b$
therefore retains that rate. The actual centered factor is
\[
 \mathcal S(k)=\exp\bigl(u^4 t_{K(k)}+5\bigr),
\]
where $t_{K(k)}\leq B_{K(k)}$ eventually. Its logarithm divided
by $t_n$ is at most $\epsilon s_{K(k)}/t_n+5/t_n$.
Taking the upper limit and then $\epsilon\downarrow0$ proves its
upper rate zero.

\subsection{Keeping the match in the continuum limit}
For fixed intermediate levels and a fixed cluster on the $j$-grid,
require
\[
 8\max_r|x_r|<a_{m+1},\qquad
 |x_r-x_s|>2a_j B_j^2\quad(r\ne s).
\]
Cross-cell separation persists during every intervening descent.
The initial room condition is
$(B_{m+1}^2+1)a_{m+1}\leq a_k$, which follows eventually from
subpower $B_{m+1}$ and $a>b$. At an intermediate level $i\geq m+1$,
put $L=a_{m+1}/a_i$. Distinct parent cell centers at level $m+1$
have lattice separation at least $a_k/a_{m+1}\geq B_{m+1}^2+1$.
Their embeddings at level $i$ therefore have separation at least
$L(B_{m+1}^2+1)$. The children each deviate from their embedded
parent by less than $L/2$. The triangle inequality leaves separation
greater than $LB_{m+1}^2$. The purchased support propagation bound
$B_i^2\leq LB_{m+1}^2$ then gives the required separation at level
$i$. This propagation bound is one of the local RG inputs retained
above. The argument applies simultaneously to every pair of marks
belonging to different copied cells.

For sufficiently large periods there is no wrapping in this finite
configuration. Its points have the required periodic separation and
have common parent zero at level $m+1$: each is inside half the
parent block by the first bound and the maximum-coordinate estimate.
For sufficiently fine cutoff $N$, the fixed stencils of the cluster
fields and of every degree-$q$ field fit in the centered positive
chamber. The cutoff comparison can consequently be applied on the
ordered tail, first in period and then in $N$.

Choose its labels $Q_{N,p}$ on that tail and set them to zero
elsewhere. Their quotas are bounded by the same $V$. Equation
\eqref{eq:ope-whole-distance} bounds their reference norms by a
constant depending on the fixed matching level, so finite-dimensional
compactness gives a convergent family on a finer nontrivial filter.
Every eventual discrepancy and quota bound persists on that same
filter. Section~\ref{sec:ope-phase-prices} converts the discrepancy
bound to all character phases, and Section~\ref{sec:ope-filtration}
gives the physical low-degree distances and the whole-label bound.
These are all statements about the same limiting label $Q$.

Finally the geometric requirements hold for the grid clusters of
Section~\ref{sec:ope-grid}. If their scale is $x$ with $a<x<c$,
their radius has rate $-x<-a$, so they fit strictly inside the
$m+1$ block. Their pairwise separation is eventually larger than
$e^{-(x+\epsilon)t_n}$, while $a_jB_j^2$ has rate $-c$ because
$B_j$ is subpower. Choosing $x+\epsilon<c$ proves the separation
condition. The fixed support-radius room at the centered scale is
also automatic: a fixed radius plus a subpower truncation radius,
multiplied by $a_m$, has rate $-a<-b$. It therefore fits inside
the centered scale $a_k$. This completes the accounting and extraction
of the local matching input, conditional on the explicit analytic
representation and comparison estimates above.
